\section{Limitaciones y amenazas a la validez}
\label{sec:limitaciones}

Los resultados deben interpretarse dentro del alcance del modelo, de las
configuraciones reducidas analizadas y del nivel de certificación alcanzado
en cada experimento. Esta sección resume las principales limitaciones
criptográficas, matemáticas y computacionales del trabajo.


\subsection{Alcance de las configuraciones reducidas}
\label{subsec:alcance_configuraciones_reducidas}

El estudio considera tamaños de palabra

\[
z \in \{4,8\},
\]

correspondientes a estados de 100 y 200 bits. Estas configuraciones
conservan la organización tridimensional de Keccak y la secuencia de
transformaciones

\[
\theta,
\quad
\rho,
\quad
\pi,
\quad
\chi,
\quad
\iota,
\]

pero no corresponden a la instancia Keccak-$f[1600]$ utilizada por SHA-3.

Por tanto, los resultados no deben interpretarse como una evaluación
directa de la seguridad de SHA-3 ni como mínimos aplicables
automáticamente a tamaños de palabra mayores. Su finalidad es estudiar la
formulación MILP y la propagación de diferencias en configuraciones
computacionalmente manejables.

De forma similar, el análisis se limita a un máximo de tres rondas. El
crecimiento de las variables binarias, las restricciones auxiliares y el
espacio de ramificación dificulta extender directamente la formulación a
más rondas mediante CBC.


\subsection{Alcance de la métrica de actividad}
\label{subsec:alcance_metrica_actividad}

La variable principal es el número de aplicaciones locales de $\chi$ que
reciben una diferencia de entrada no nula. Esta medida permite cuantificar
la cantidad de componentes no lineales atravesados por una trayectoria,
pero no representa por sí sola su probabilidad diferencial.

Dos trayectorias con el mismo número de S-boxes activas pueden contener
transiciones locales diferentes y, en consecuencia, presentar
probabilidades distintas. El modelo tampoco calcula directamente:

\begin{itemize}
    \item el peso diferencial de cada transición de $\chi$;
    \item la probabilidad acumulada de una trayectoria;
    \item el número de características compatibles con un mismo soporte;
    \item correlaciones lineales;
    \item la complejidad de un ataque contra SHA-3.
\end{itemize}

En consecuencia, la minimización de actividad identifica trayectorias
estructuralmente poco dispersas, pero no garantiza que sean las
trayectorias diferenciales más probables.

Una extensión natural consistiría en incorporar pesos asociados a las
transiciones diferenciales de $\chi$ y minimizar el peso total, en lugar
de utilizar únicamente indicadores binarios de actividad.


\subsection{Brecha de certificación en tres rondas}
\label{subsec:limitacion_tres_rondas}

El nivel de certificación no es igual para todos los experimentos.

Para una y dos rondas se obtuvo una coincidencia entre las cotas inferiores
y los testigos factibles. Por tanto, los valores reportados corresponden a
mínimos globales.

Para tres rondas solo se dispone de los intervalos

\[
5
\leq
N_{\mathrm{act}}^{\min}(4,3)
\leq
13,
\]

y

\[
5
\leq
N_{\mathrm{act}}^{\min}(8,3)
\leq
14.
\]

Las cotas inferiores son globales y los extremos superiores están
respaldados por testigos concretos. Sin embargo, no se demostró que no
existan trayectorias con actividades comprendidas dentro de estos
intervalos.

Por esta razón, los valores 13 y 14 no constituyen mínimos globales
certificados.


\subsubsection{Restricción de la familia \texorpdfstring{$2+2+c$}{2+2+c}}

La búsqueda de tres rondas se restringió a trayectorias cuya actividad en
las dos primeras rondas es

\[
2+2.
\]

Dentro de esta familia, la enumeración fue exhaustiva y produjo los mejores
valores

\[
2+2+9=13
\]

para $z=4$, y

\[
2+2+10=14
\]

para $z=8$.

No obstante, una trayectoria globalmente mejor podría presentar una
distribución diferente, por ejemplo,

\[
1+a+b,
\qquad
3+a+b,
\qquad
m_0+m_1+m_2,
\]

siempre que satisfaga las restricciones exactas de propagación.

En consecuencia, la afirmación correcta es:

\begin{quote}
La búsqueda es exhaustiva dentro de la familia $2+2+c$, pero no dentro
del espacio completo de trayectorias de tres rondas.
\end{quote}

Cerrar los intervalos requeriría explorar otras distribuciones de actividad,
fortalecer la formulación MILP o emplear métodos de descomposición más
generales.


\subsection{Dependencia del solver y del entorno}
\label{subsec:dependencia_solver_entorno}

CBC permitió resolver los casos de una y dos rondas, validar testigos y
demostrar la infactibilidad de determinados problemas auxiliares. Sin
embargo, su capacidad de certificación disminuye al aumentar el tamaño del
modelo.

El crecimiento de $z$ y $R$ incrementa:

\begin{itemize}
    \item las variables de los estados izquierdo y derecho;
    \item las variables de diferencia;
    \item las variables auxiliares para XOR y productos;
    \item las restricciones de cada transformación;
    \item el espacio explorado mediante ramificación y acotación.
\end{itemize}

Una terminación por límite de tiempo no demuestra que el modelo sea
infactible ni que la mejor solución encontrada sea óptima. En este trabajo
solo se utilizan como certificadas las conclusiones respaldadas por un
estado concluyente del solver o por un testigo reconstruido
independientemente.

El uso de un solver comercial, como Gurobi, podría reducir los tiempos de
resolución mediante mejores procedimientos de \emph{presolve}, cortes,
heurísticas y paralelización. No obstante, un solver más rápido no elimina
la necesidad de distinguir entre:

\begin{itemize}
    \item una solución factible;
    \item una cota inferior;
    \item una prueba de infactibilidad;
    \item una prueba de optimalidad global.
\end{itemize}

Los tiempos de ejecución también dependen del procesador, la memoria, las
versiones de Python, PuLP y CBC, y la carga concurrente del sistema. Por
ello, se reportan como referencias de reproducibilidad y no como una
comparación formal de rendimiento.


\subsection{Validez de la implementación}
\label{subsec:validez_implementacion}

La validez interna se fortalece mediante:

\begin{itemize}
    \item una implementación funcional independiente;
    \item comparación bit a bit entre la referencia y el modelo MILP;
    \item reconstrucción de estados, diferencias y soportes;
    \item generación automática de resultados;
    \item pruebas unitarias y de integración;
    \item control de versiones y testigos almacenados.
\end{itemize}

Estas medidas reducen el riesgo de aceptar resultados producidos por
errores de programación o de extracción de variables. Sin embargo, como en
cualquier implementación computacional, las pruebas automatizadas no
constituyen una demostración formal de la corrección total del software.

Los testigos almacenados demuestran la existencia de trayectorias con una
actividad determinada, pero no describen todas las soluciones con el mismo
valor. Pueden existir diferentes soportes, estados absolutos o soluciones
relacionadas mediante simetrías que produzcan la misma actividad.

Por tanto, los testigos deben interpretarse como certificados de
factibilidad y no como representantes únicos de todas las trayectorias
óptimas o cuasi-óptimas.


\subsection{Correspondencia entre intentos y rondas}
\label{subsec:limitacion_intentos_rondas}

La relación

\[
R(t)=
\begin{cases}
1, & t<10,\\
2, & 10\leq t<20,\\
3, & 20\leq t<30
\end{cases}
\]

se incorporó para responder al planteamiento de la actividad académica.

Esta relación no pertenece a la definición de Keccak y no representa una
conexión criptográfica demostrada entre el número de intentos de un
atacante y el número de rondas de la permutación. En la implementación,
los intentos funcionan únicamente como un parámetro externo para seleccionar
uno de los tres escenarios experimentales.

Por esta razón, las conclusiones se expresan principalmente en términos de
$z$ y $R$, y no como una relación general entre intentos y seguridad.


\subsection{Síntesis de las amenazas a la validez}
\label{subsec:sintesis_amenazas}

La Tabla~\ref{tab:amenazas_validez} resume las principales amenazas y las
medidas adoptadas para reducirlas.

\begin{table}[htbp]
    \centering
    \caption{Principales amenazas a la validez y medidas de mitigación.}
    \label{tab:amenazas_validez}
    \begin{tabular}{
        p{0.24\textwidth}
        p{0.31\textwidth}
        p{0.35\textwidth}
    }
        \hline
        \textbf{Amenaza}
        &
        \textbf{Consecuencia}
        &
        \textbf{Mitigación}
        \\
        \hline

        Tamaños y rondas reducidos
        &
        Los resultados no se trasladan automáticamente a Keccak-$f[1600]$.
        &
        Delimitación explícita del alcance y uso de configuraciones
        estructuralmente válidas.
        \\[4pt]

        Métrica basada en actividad
        &
        No se obtiene la probabilidad diferencial completa.
        &
        Interpretación estructural y propuesta de incorporar pesos
        diferenciales.
        \\[4pt]

        Brecha en tres rondas
        &
        No se conoce el mínimo global exacto.
        &
        Reporte mediante intervalos y distinción entre cotas y óptimos.
        \\[4pt]

        Búsqueda $2+2+c$
        &
        Pueden existir mejores distribuciones fuera de la familia.
        &
        Declaración explícita de la restricción y conservación de los
        testigos como cotas superiores globales.
        \\[4pt]

        Dependencia del solver
        &
        Los tiempos y la capacidad de certificación pueden variar.
        &
        Uso exclusivo de estados certificados y validación independiente.
        \\[4pt]

        Errores de implementación
        &
        Una formulación incorrecta podría producir resultados inválidos.
        &
        Implementación de referencia, pruebas automatizadas y
        reconstrucción de testigos.
        \\
        \hline
    \end{tabular}
\end{table}

Estas limitaciones no invalidan los resultados obtenidos. Definen el
contexto dentro del cual pueden considerarse correctos y evitan extender
las conclusiones más allá de la evidencia disponible.